\chapter{Fazit}
\label{chap:fazit}

Ziel dieser Arbeit war es, stabile und kontrollierbare Lernprozesse für Reinforcement-Learning-Agenten in komplexen Survival-Szenarien zu ermöglichen. Erste Experimente im Multi-Agenten-Setting zeigten, dass ohne klar isolierte Lernschritte keine verlässliche Konvergenz der Policy erreicht werden konnte. Aus diesem Grund wurde ein modularer Ansatz verfolgt, bei dem zentrale Teilstrategien wie Navigation, Exploration und Flucht zunächst separat trainiert und anschließend im Rahmen eines Curriculum-Learning-Verfahrens in ein umfassendes Survival-Szenario integriert wurden.

Die Ergebnisse zeigen, dass mehrere Faktoren entscheidend zur Stabilität des Trainings beitragen. Dazu gehören eine eng abgestimmte Hyperparameterkonfiguration, klar strukturierte Rewardfunktionen sowie eine gezielte Gestaltung der Umgebungen. Besonders das Reward-Shaping erwies sich als wesentlich, um zielgerichtetes Verhalten zu fördern. Zu aggressive oder widersprüchliche Belohnungskomponenten führten häufig zu instabilen Strategien oder einem vollständigen Zusammenbruch des Lernprozesses. Darüber hinaus hatte die Struktur des Beobachtungs- und Aktionsraums einen deutlichen Einfluss auf das Lernverhalten. Die Verwendung eines diskreten Aktionsraums in Kombination mit vereinfachten Richtungsvektoren erleichterte die Entwicklung effektiver Strategien erheblich.

Im direkten Vergleich mit Modellen, die ohne Vortraining trainiert wurden, zeigten vortrainierte Modelle im Survival-Szenario eine ähnliche Gesamtleistung. Gleichzeitig entwickelten diese Modelle spezifische strategische Präferenzen. So erreichten Navigationsmodelle bereits früh erfolgreiche Episoden, Explorationsmodelle wiesen eine insgesamt höhere Erfolgsquote auf, während Fluchtmodelle Schwierigkeiten hatten, ihr Verhalten an das komplexere Szenario mit mehreren Zielen anzupassen. Diese Muster fanden sich auch in der Analyse des Fähigkeitserhalts wieder: Die Fähigkeit zur Exploration konnte durch das anschließende Fine-Tuning sogar weiter verbessert werden, das Navigationsverhalten blieb weitgehend stabil, wohingegen die Fluchtfähigkeit deutlich zurückging.

Die Ergebnisse verdeutlichen, dass modulare Trainings- und Transferstrategien das Lernverhalten gezielt beeinflussen können. Sie ermöglichen einerseits eine beschleunigte Entwicklung gewünschter Strategien, können andererseits aber auch bestehende Verhaltensmuster verfestigen, die im neuen Anwendungskontext nicht mehr optimal sind. Die entwickelte Simulationsumgebung stellte hierbei einen wesentlichen Vorteil dar, da sie eine systematische und reproduzierbare Analyse einzelner Teilfähigkeiten sowie deren Zusammenspiel im Gesamtsystem ermöglichte.

Die Aussagekraft der vorliegenden Ergebnisse ist begrenzt durch die Beschränkung auf eine einzige Modellarchitektur in Form eines LSTM-basierten PPO-Ansatzes sowie durch die Anzahl der durchgeführten Trainingsläufe. Zukünftige Arbeiten könnten hiervon ausgehend alternative Modellvarianten untersuchen, beispielsweise attention-basierte Netzwerke oder hierarchisch strukturierte Steuermechanismen.

Zusammenfassend lässt sich festhalten, dass sorgfältig entworfene Teilaufgaben und ein kontrollierter Transfermechanismus eine tragfähige Grundlage darstellen, um das Strategielernen in komplexen Reinforcement-Learning-Umgebungen robuster und zielgerichteter zu gestalten. Entscheidend ist dabei, dass der Wissenstransfer bewusst gestaltet und an die Anforderungen des Zielkontexts angepasst wird. Die zu Beginn dieser Arbeit formulierten Ziele wurden vollständig erreicht. Dazu zählen die Untersuchung modularer Lernverfahren, die Analyse von Transferprozessen und des Fähigkeitserhalts sowie die Entwicklung einer kontrollierbaren Simulationsumgebung.
